\chapter{拓扑空间.}

极限和连续的概念可以追溯到古代；若不从这个观点出发系统地研究不仅数学家们、而且还有希腊哲学家们特别是亚里士多德，若不追随这些思想在文艺复兴时期数学以及微分积分学开端中的演变，便不可能写出一部完整的历史。这样一项研究做起来固然非常有意思，但它会远远超出本篇笔记的框架。

应当把拓扑学——如同现代数学的许多其他分支一样——的创立者视为黎曼：事实上正是他第一个力图提炼出拓扑空间的概念，构想出关于这些空间的一门自治理论的思想，定义了在拓扑学日后发展中起最大作用的不变量（“贝蒂数”），并给出了它对分析的首批应用（阿贝尔积分的周期）。但 19 世纪上半叶的思想运动并非没有在不止一个方面为黎曼做好准备。的确，把数学安放在坚实基础之上的愿望——它是整个 19 世纪直至我们今天如此众多重要研究的起因——已经导致了收敛级数的概念和趋于一个极限的数列的概念（柯西、阿贝尔）以及连续函数的概念（波尔查诺、柯西）的正确定义。另一方面，复数——或者像当时人们所说的那样，“虚数”（在 18 世纪有时还被冠以“不可能的数”之名）——的几何表示（用平面上的点），这种归功于阿尔冈和高斯（见第 161 页）的表示法，已为大多数数学家所熟悉：它所构成的进步，与我们今天在希尔伯特空间研究中采用几何语言属于同一层次，并且以胚芽的形态包含着对每一个可连续变化的对象作几何表示的可能性；高斯本来就因其关于几何基础、非欧几何、曲面的研究而很自然地被引向这类观念，他看来已经预见到了这种可能性，因为他在（独立于阿尔冈和法国数学家们）定义虚数的几何表示时使用了“二重广延的大小”一语（{[}124 a{]}，第 II 卷，第 101–103 页及第 175–178 页）。

一方面是他关于代数函数及其积分的研究，另一方面是他关于几何基础的思考（在很大程度上得自他对高斯著作的研读），正是这两者引导黎曼表述出一个恰好就是现代拓扑学研究纲领的纲领，并使这一纲领的实现有了开端。例如，他在其《阿贝尔函数论》（{[}259 a{]}，第 91 页）中是这样说的：

“在研究通过积分恰当微分而得到的函数时，位置分析的几条定理几乎是不可缺少的。这个名称曾为莱布尼茨所用，尽管含义也许有些不同，我们要把它理解为连续量理论中研究这些量的那一部分：不是把这些量看作独立于其位置、可用彼此量度的东西来研究，而是抽掉一切度量的观念，只研究它们的位置关系和包含关系。我保留日后以一种完全独立于一切度量的方式处理这个对象的权利 \ldots{}”。

而在他著名的就职演讲《论作为几何基础的假设》（{[}259 a{]}，第 272 页）中：

“\ldots 多次广延的一般大小概念，\(^{1}\)它把空间大小作为特殊情形包含在内，至今完全没有得到研究 \ldots{}（第 272 页）”。

“\ldots 大小的概念预设了一个可取各种不同规定的元素。根据能否从一个规定连续地过渡到另一个规定，这些规定构成一个连续流形（它们将被称为该流形的点）或一个离散流形（第 273 页）”。

“\ldots 度量在于把要比较的各个量相互叠置；因此，要进行度量，就必须有一种把一个量移置到另一个量上的手段。在缺乏这种手段时，只有当一个量构成另一个量的一部分，才可能比较这两个量\ldots{}就此所能作的研究，构成了大小理论中独立于度量的那一部分，在这里，量既不被看作独立于其位置而存在，也不被看作可用一个度量单位来表达，而被看作一个流形的各个部分。这样的研究在数学的许多部分中已成为必需，对解析多值函数的理论来说尤其如此 \ldots{}（第 274 页）”。

“\ldots 在一个给定的流形中，位置的确定，凡属可能，都归结为有限多个数值规定。诚然，存在这样一些流形，其中位置的确定所需要的不是有限多个规定，而是一个无穷序列或一个连续的量规定之流形。这样的流形例如由一个函数在一个给定区域中的各种可能规定、由空间中一个图形的位置等等构成”（第 276 页）。

在这最后一句里，可以注意到函数空间研究的最初想法；此外，早在黎曼的博士论文中，同样的想法就已经得到表达：他在谈到以狄利克雷原理之名著称的最小值问题时说，“这些函数的全体构成一个连通的、自身封闭的区域”（{[}259 a{]}，第 30 页），这虽然形式尚不完善，却已经是希尔伯特日后给出的狄利克雷原理证明、乃至函数空间对变分法的大多数应用的胚胎。

如前所述，黎曼为这一宏伟纲领的实现开了头：他定义了“贝蒂数”，先是一个曲面的贝蒂数（{[}259 a{]}，第 92–93 页），然后（同前引文，第 479–482 页；另参见{[}259 c{]}）是任意维数的一个流形的贝蒂数，并把这一定义应用于积分理论；这些结果开创了代数拓扑学，这门数学分支自 20 世纪初以来发展不断加速，其历史不应在此追述。

至于拓扑空间的一般理论，即黎曼所预见的那种理论，它要得到发展，就必须先对实数理论、数集理论、直线上的点集理论以及平面和空间中的点集理论，作比黎曼时代更为系统的研究：这一研究一方面联系着关于无理数本性的探讨（波尔查诺的偏于哲学，戴德金的则本质上是数学的），另一方面联系着实变量函数论的进展（黎曼本人以他的积分定义和他的三角级数理论为它作出了重要贡献，杜布瓦-雷蒙、迪尼、魏尔斯特拉斯等人的工作也以它为对象）；这是 19 世纪下半叶的工作，而其中最关键的是康托尔，他第一个定义了（先是在直线上，然后在 n 维欧几里得空间中）聚点、闭集和完全集的概念，并得到了关于这些集在直线上的结构的本质结果（参见第 155 页）：对此不仅应查阅康托尔的著作{[}47{]}，还应查阅他与戴德金的那批非常有趣的通信{[}48{]}，在那里还可以清楚地看到把维数视为拓扑不变量的思想。

理论的后续进展，例如在舍恩弗利斯的书{[}275 a{]}中以半历史、半系统的形式作了陈述：其中远远最为重要的是博雷尔–勒贝格定理（最先由博雷尔就直线上的一个闭区间和覆盖它的可数开区间族证明）。

康托尔的思想起初遇到了相当激烈的反对（参见第 28 页）。但至少他关于直线上和平面上的点集的理论，很快就被法国和德国的函数论学派（若尔当、庞加莱、克莱因、米塔-列夫勒，随后是阿达马、博雷尔、贝尔、勒贝格等）所使用和广泛传播：特别是博雷尔丛书的开头几卷，每一卷都含有这一理论的一个初等陈述（例如见{[}32 a{]}）。随着这些思想的传播，人们开始以各种方式思考把它们应用于不再是点集、而是曲线之集或函数之集的可能性：例如早在 1884 年，这一想法就已见诸阿斯科利一篇论文的标题“论曲线族的极限曲线”{[}10{]}，并且它还表达在阿达马 1896 年向苏黎世数学家大会所作的一次通报中{[}141{]}；它也与沃尔泰拉 1887 年引入“线函数”以及创立“泛函演算”（即变元为函数的函数论）紧密相联（关于这一点，可参看沃尔泰拉关于泛函分析的著作{[}322 b{]}）。另一方面，在那篇著名的论文（{[}163 a{]}，第 III 卷，第 10–37 页）中，希尔伯特在这一点上重拾黎曼的思想，证明了狄利克雷原理中最小值的存在，并开创了变分法的“直接方法”，在其中可以清楚地看到考虑使波尔查诺–魏尔斯特拉斯原理成立的函数集（也就是说，每个序列都含有一个收敛的子序列）有何意义；这样的集合不久确实发挥了重要作用，不仅在变分法中，而且在实变量函数论（阿斯科利、阿尔泽拉）和复变量函数论（维塔利、卡拉泰奥多里、蒙泰尔）中都是如此。最后，函数方程的研究，尤其是弗雷德霍姆对以他的名字命名的那类方程的求解{[}116{]}，使人们熟悉了把函数看作变元、把函数之集看作点集之类似的考虑，对这样的集合使用几何语言，与对 n 维欧几里得空间的点使用几何语言同样自然（这种空间同样也逃出了“直观”，并因此而在很长时间里始终是许多数学家不信任的对象）。特别地，希尔伯特关于积分方程的那些值得纪念的工作{[}163 b{]}，以埃哈德·施密特{[}274 b{]}对希尔伯特空间的定义和几何研究而告终，它与欧几里得几何完全类似（见第 212 页）。

与此同时，公理化理论的概念已变得愈来愈重要，这首先要归功于关于几何基础的大量著作，其中希尔伯特的{[}163 c{]}产生了尤其决定性的影响；正是在这项工作期间，希尔伯特早在 1902 年（{[}163 c{]}，第 180 页）就被引导去精确陈述黎曼意义上的“二重广延流形”的第一个公理化定义，照他的说法，这个定义构成“对位置分析作严格公理化处理的基础”，而且已经使用了邻域（其含义受到希尔伯特当时所限定的问题之要求的限制）。

把点集的性质与函数集的性质中共同的东西提炼出来（而不引入“距离”的概念）的最初尝试，是由弗雷歇{[}115 a{]}和 F. 里斯{[}260 b{]}作出的；但前者从可数极限的概念出发，未能为不可度量化的空间构造出一个合用而富有成果的公理体系；至少他认识到了波尔查诺–魏尔斯特拉斯原理与博雷尔–勒贝格定理之间的关系；正是在这一关联中他引入了“紧”这个词，尽管其含义与今天赋予它的含义略有不同。至于 F. 里斯，他从聚点（或者更确切地说，从与之等价的“导集”）的概念出发，他的理论仍然是不完全的，况且始终只是一个概要。

随着豪斯多夫（{[}152 a{]}，第 7-8-9 章），今天所理解的一般拓扑学开始了。他重新拾起邻域的概念，并且懂得从希尔伯特关于平面上邻域的公理中挑选出那些能够同时赋予他的理论以一切想要的精确性和一般性的公理。他推演这些公理的推论的那一章，至今仍是公理化理论的一个典范：抽象，然而从一开始就适用于应用。很自然，这是一般拓扑学后续研究的出发点，尤其是莫斯科学派的工作，其大部分指向度量化问题（参见第 166 页）：我们在这里特别应当从中记住亚历山德罗夫和乌雷松对紧空间的定义（以“双紧空间”之名），然后是吉洪诺夫{[}313{]}关于紧空间的乘积的紧性的证明。最后是 H. 嘉当{[}53{]}引入的滤子，它为各种各样的应用提供了一个非常宝贵的工具（在其中它是“摩尔–史密斯收敛”概念{[}227{]}的一个有利的替代品），并且借助关于超滤子的定理，最终实现了理论的澄清与简化。

